\chapter{Verdelen en delen}\label{ch:g3:division}

Twintig knikkers eerlijk verdeeld over vier kinderen: hoeveel krijgt
ieder? Delen antwoordt op alle verdeelvragen --- en ook op de
groepeervragen --- en het is de tweeling van het vermenigvuldigen: het
ene maakt het andere ongedaan.

\section{Twee kanten van het delen}

\begin{definition}[Delen]\label{def:g3:division:def}
De deling $20 \div 4$ antwoordt op twee soorten vragen:
\begin{itemize}
  \item \emph{verdelen}: $20$ knikkers verdeeld over $4$ kinderen ---
        hoeveel knikkers krijgt \emph{ieder}? ($5$ elk);
  \item \emph{groeperen}: $20$ knikkers in zakjes van $4$ --- hoeveel
        \emph{zakjes}? ($5$ zakjes).
\end{itemize}
Beide vragen hebben dezelfde uitkomst, want beide maken dezelfde
vermenigvuldiging ongedaan: $4 \times 5 = 20$.
\end{definition}

\begin{omfigure}
\begin{tikzpicture}[scale=0.55]
  \foreach \g in {0,1,2,3} {
    \draw[omExo, rounded corners] (\g*4.9-0.65,-0.7) rectangle
      (\g*4.9+3.85,0.7);
    \foreach \c in {0,...,4} \fill[omDef] (\g*4.9+\c*0.8,0) circle (0.27);
  }
  \node[below, font=\small] at (8.95,-1.15)
    {$20$ knikkers, $4$ gelijke groepjes: $5$ in elk groepje};
\end{tikzpicture}

{\small $20$ verdeeld over $4$ gelijke groepjes. Groeperen per $4$ zou
$5$ groepjes geven --- dezelfde deling.}
\end{omfigure}

\begin{method}[Delen met de tafels]\label{met:g3:division:tables}
Om $42 \div 6$ uit te rekenen, vraag je: ``$6$ keer \emph{wat} is
$42$?'' De tafel van $6$ antwoordt: $6 \times 7 = 42$, dus
$42 \div 6 = 7$. Elke deelvraag is een vermenigvuldigvraag, achterstevoren
gelezen.
\end{method}

\section{Als het niet precies uitkomt}

\begin{definition}[Rest]\label{def:g3:division:remainder}
$23$ knikkers verdelen over $4$ kinderen geeft $5$ knikkers elk --- en
$3$ knikkers die \emph{overblijven}, te weinig voor nog een ronde. We
schrijven
\[
23 = 4 \times 5 + 3 ,
\]
en zeggen: het \emph{quotiënt}\index{quotiënt} is $5$, de
\emph{rest}\index{rest} is $3$. De rest is altijd kleiner dan het aantal
kinderen --- anders zou er nog een ronde verdeeld kunnen worden.
\end{definition}

\begin{omfigure}
\begin{tikzpicture}[scale=0.55]
  \foreach \g in {0,1,2,3} {
    \draw[omExo, rounded corners] (\g*4.9-0.65,-0.7) rectangle
      (\g*4.9+3.85,0.7);
    \foreach \c in {0,...,4} \fill[omDef] (\g*4.9+\c*0.8,0) circle (0.27);
  }
  \foreach \c in {0,1,2} \fill[omThm] (19.7+\c*0.8,0) circle (0.27);
  \node[above, font=\footnotesize, omThm] at (20.5,0.5) {blijft over};
  \node[below, font=\small] at (10.2,-1.15)
    {$23 = 4 \times 5 + 3$: quotiënt $5$, rest $3$ (in rood)};
\end{tikzpicture}

{\small $23$ verdelen over $4$: na $5$ ronden blijven er $3$ knikkers
over --- te weinig voor een $6$de ronde.}
\end{omfigure}

\begin{example}\label{ex:g3:division:remainder}
Deel $58$ door $7$. Loop de tafel van $7$ langs: $7 \times 8 = 56$ is het
grootste \omterm{def:g2:mult:def}{product} dat onder $58$ blijft ($7 \times 9 = 63$ is te veel).
Dus
\[
58 = 7 \times 8 + 2 :
\]
\omterm{def:g3:division:remainder}{quotiënt} $8$, \omterm{def:g3:division:remainder}{rest} $2$. Controle: $56 + 2 = 58$, en $2 < 7$.
\end{example}

\begin{method}[De rest uitleggen]\label{met:g3:division:interpret}
Ga na het delen terug naar het verhaal:
\begin{enumerate}
  \item ``hoeveel elk / hoeveel volle groepjes?'' --- het \omterm{def:g3:division:remainder}{quotiënt};
  \item ``hoeveel blijft er over?'' --- de \omterm{def:g3:division:remainder}{rest};
  \item ``hoeveel dozen om alles in te doen?'' --- het \omterm{def:g3:division:remainder}{quotiënt}
        \emph{plus één} als de \omterm{def:g3:division:remainder}{rest} niet \omterm{def:g1:counting:zero}{nul} is.
\end{enumerate}
\end{method}

\begin{example}\label{ex:g3:division:interpret}
$58$ leerlingen gaan kanoën, $7$ per kano. $58 = 7 \times 8 + 2$: acht
kano's zijn vol, en er blijven $2$ leerlingen over --- die hebben ook een
kano nodig. Er moeten dus $9$ kano's worden gereserveerd.
\end{example}

\section{Helften, derde delen, kwarten}

\begin{definition}[Helft, derde deel, kwart]\label{def:g3:division:half}
De \emph{helft}\index{helft} van een getal is de uitkomst van de deling
door $2$; het \emph{derde deel} van de deling door $3$; het
\emph{kwart}\index{kwart} van de deling door $4$. Zo is de helft van
$18$ gelijk aan $9$, het derde deel van $18$ gelijk aan $6$, en het kwart
van $36$ gelijk aan $9$.
\end{definition}

\begin{example}[Dubbels en helften horen bij elkaar]\label{ex:g3:division:double}
De \omterm{def:g3:division:half}{helft} van $46$: omdat $46 = 40 + 6$, geeft de \omterm{def:g3:division:half}{helft} van elk deel
$20 + 3 = 23$. Controleer met het dubbele: $23 + 23 = 46$. Wie de dubbels
uit het hoofd kent ($25 \to 50$, $45 \to 90$, $150 \to 300$), vindt
\omterm{def:g3:division:half}{helften} in een oogwenk.
\end{example}

\section{Oefeningen}

\begin{exercise}[$\star$]\label{exo:g3:division:1}
Reken uit met de tafels achterstevoren: $36 \div 4$; \quad $63 \div 9$;
\quad $40 \div 5$; \quad $56 \div 8$.
\end{exercise}

\begin{exercise}[$\star$]\label{exo:g3:division:2}
Schrijf bij elke deling de gelijkheid $a = b \times q + r$ op en geef het
\omterm{def:g3:division:remainder}{quotiënt} en de \omterm{def:g3:division:remainder}{rest}: $17 \div 5$; \quad $30 \div 4$; \quad
$50 \div 6$.
\end{exercise}

\begin{exercise}[$\star$]\label{exo:g3:division:3}
$28$ knikkers worden verdeeld over $6$ kinderen. Hoeveel knikkers krijgt
ieder, en hoeveel blijven er over? Teken de verdeling als dat helpt.
\end{exercise}

\begin{exercise}[$\star$]\label{exo:g3:division:4}
Reken uit: de \omterm{def:g3:division:half}{helft} van $68$; het derde deel van $27$; het \omterm{def:g3:division:half}{kwart} van
$48$; de \omterm{def:g3:division:half}{helft} van $150$.
\end{exercise}

\begin{exercise}[$\star$]\label{exo:g3:division:5}
Waar of niet waar? Controleer met een vermenigvuldiging.
``$45 \div 5 = 9$''; \quad ``$52 \div 6 = 8$ \omterm{def:g3:division:remainder}{rest} $4$'';
\quad ``$70 \div 8 = 9$ \omterm{def:g3:division:remainder}{rest} $2$''.
\end{exercise}

\begin{exercise}[$\star$]\label{exo:g3:division:6}
$32$ foto's worden in een album geplakt, $4$ per bladzijde. Hoeveel
bladzijden zijn er nodig? Welke kant van het delen is dit (verdelen of
groeperen)?
\end{exercise}

\begin{exercise}[$\star$]\label{exo:g3:division:7}
Een touw van $75$ cm wordt in stukken van $9$ cm geknipt. Hoeveel hele
stukken krijg je, en hoe lang is het stukje dat overblijft?
\end{exercise}

\begin{exercise}[$\star$]\label{exo:g3:division:8}
$50$ kinderen gaan op excursie in auto's met $4$ plaatsen. Hoeveel auto's
zijn er nodig? (Pas op: iedereen moet mee!)
\end{exercise}

\begin{exercise}[$\star$]\label{exo:g3:division:9}
Zoek de ontbrekende getallen: $? \div 7 = 6$; \quad
$54 \div ? = 6$; \quad de \omterm{def:g3:division:half}{helft} van $?$ is $35$.
\end{exercise}

\begin{exercise}[$\star\star$]\label{exo:g3:division:10}
Amina verdeelt haar $47$ stickers zo eerlijk mogelijk onder haar $5$
vriendinnen en houdt wat overblijft zelf. Hoeveel stickers krijgt elke
vriendin, en hoeveel houdt Amina?
\end{exercise}

\begin{exercise}[$\star\star$]\label{exo:g3:division:11}
Welke \omterm{def:g3:division:remainder}{resten} zijn er allemaal mogelijk bij een deling door $6$? Wat is
het grootste getal waarvan de deling door $6$ \omterm{def:g3:division:remainder}{quotiënt} $7$ heeft?
\end{exercise}
